\section{问题重述与分析}
\subsection{问题重述}
热风烘干通过烘房与药材表面的热量、水分交换改变药材内部状态。题给药材近似为长25 cm、初始半径2 cm的圆柱，初始温度28 $^\circ$C、干基含水率2.55 kg/kg。附件1记录初期烘房温湿环境，附件2记录烘干中的半径变化。本文围绕以下四个任务建立模型。

问题一采用附录2参数，计算前1800 s的径向温度与含水率分布，在指定七个时刻、五个半径位置列出结果，并按1 s、0.1 cm间隔输出完整数据。问题二从初始状态统一采用附录3经验物性，计算前3 h的温湿分布。问题三在问题二模型基础上确定药材各处含水率低于0.15 kg/kg的累计时间，并输出每6 h的径向含水率。问题四利用附件2的半径轨迹与附录4物性，重新计算收缩条件下的烘干时间和含水率分布。

问题一、二的细时间结果分别保存至\texttt{result1.xlsx}、\texttt{result2.xlsx}；问题三、四按60 s输出至\texttt{result3.xlsx}、\texttt{result4.xlsx}。所有结果表保留四位小数，数值推进和达标判定使用未舍入值。

\subsection{模型的递进关系}
四问的共同核心是内部扩散与表面交换的共同作用。集中参数模型只能描述平均状态，难以直接满足题目对不同径向位置的要求。Hussain和Dincer研究了圆柱湿物体干燥中的二维热湿传递问题~\cite{hussain2003}；Kaya等讨论了圆柱湿物体强制对流干燥的数值建模~\cite{kaya2007}。本文保留径向空间分布，以题给物性和边界数据建立可计算的有效模型，不移植其他材料的经验参数。

问题一的热物性为常数，且扩散系数仅依赖含水率，故在不显式考虑潜热的条件下，温度场和水分场可以分别求解。问题二、三的热物性依赖含水率，水分扩散系数同时依赖局部温度，必须联合推进两个场。问题三的新增任务是对全域最大值进行事件定位，而不是利用平均含水率或早期失水斜率外推时长。问题四进一步改变空间域：既要描述材料运动，又要把固定材料坐标的解映射回当前物理半径。

上述任务均为初边值预测问题。烘房条件与半径轨迹在模型中属于输入，不引入题目没有要求的工艺优化目标。四问均从题给初始状态计算，第二问不拼接第一问末态，第四问也不读取第三问的最终状态。

\subsection{输入整理与时间延续}
附件1含241组观测，时间为0至14400 s，间隔60 s。前1800 s和前10800 s分别包含31组、181组观测。对各问所需观测区间采用保形分段三次Hermite插值（PCHIP）~\cite{pchip}。记相邻观测点为$(t_j,f_j)$，区间长度为$h_j$，$z=(t-t_j)/h_j$，则插值表达式为
\begin{equation}\label{eq_pchip}
\begin{split}
f(t)={}&(2z^3-3z^2+1)f_j+(-2z^3+3z^2)f_{j+1}\\
&+(z^3-2z^2+z)h_jm_j+(z^3-z^2)h_jm_{j+1},
\end{split}
\end{equation}
其中$m_j$由保形规则确定，$f$取烘房温度$T_a$或水分量$C_a$。式\eqref{eq_pchip}精确通过观测点，不删除小幅波动，也不作全局平滑。问题一的环境输入见图\ref{fig_input}。
\figone{q1_oven_input.pdf}{.76}{前30 min烘房温度和水分量随时间变化。曲线由附件1观测插值得到，输入升温不等于药材内部同步升温。}{fig_input}

附件1没有覆盖完整干燥时段。问题三、四在$t\le14400$ s时保留全部观测输入，随后采用恒温阶段延续假设
\begin{equation}\label{eq_late}
T_a(t)=50\ ^\circ\mathrm C,\qquad C_a(t)=0.05\ \mathrm{kg/kg},\qquad t>14400\ \mathrm s.
\end{equation}
式\eqref{eq_late}是基于恒温干燥情景的模型输入，不是后期实测数据。切换处保留原观测终值，允许边界输入发生小幅变化，内部温湿状态保持连续。

附件2含145组半径观测，覆盖0至72 h，间隔0.5 h，半径从2 cm降至1.198 cm。将半径换算为m后使用同类插值构造$R(t)$。仅在数据覆盖区间内计算，72 h不作为预设的烘干答案。
\FloatBarrier
